\scanpage{176}
Lorsque \deleted{$X$ est \unsure{l'on|on} est sous les conditions de} les anneaux localisés stricts de $X$ sont factoriels, on tire aisément de (5.3.)
\[
H^1(X,\underline{\unsure{\mathrm{Div}|\mathrm{Dv}|\mathcal{D}\mathrm{iv}}}_X)=0
\]
d'où \deleted{la proposition}

\uline{Proposition} 5.\unsure{4}. Supposons que l'on soit sous les conditions de 5.1. a), \inserted{$X$ étant irréductible de pt gén. $y$.} Alors les homomorphismes canoniques
\[
H^2(X,\mathrm{G}_{m\,X})\unsure{-|\to}H^2(X,\mathrm{R}_X^*)\longrightarrow H^2(y,\mathrm{G}_{m\,y})=\mathrm{Br}(y)
\]
sont injectifs.

Comme $\mathrm{Br}(X)\subset H^2(X,\mathrm{G}_{m\,X})$, il s'ensuit

\uline{Corollaire} 5.5. Sous les \unsure{mêmes} conditions,
\[
\mathrm{Br}(X)\unsure{-|\to}\mathrm{Br}(y)
\]
est injectif.

Ainsi, $\mathrm{Br}(X)$ s'identifie \inserted{alors} à un \unsure{sous-gpe|sous-groupe} du groupe de Brauer $\deleted{\unsure{\mathrm{B}}}\mathrm{Br}(K)$ du corps des fonctions $K=k(y)$ de $X$.

\uline{Théorème} 5.6. \inserted{Soit $X$ un préschéma, $\unsure{\zeta|\xi}\in H^2(X,\mathrm{G}_m)$} a) Si $X$ est un schéma local, et \deleted{$\zeta\in H^2(X,\mathrm{G}_{m\,X})$}, pour que $\zeta$ \deleted{\unsure{provient|provienne} d'un élément de} \inserted{soit dans $\mathrm{Br}(X)$, il} faut et s. qu'il existe un \deleted{extension} morphisme fini étale surjectif (galoisien si on veut) $f\colon X'\to X$, \deleted{tel que $\zeta|X'$ soit} qui \deleted{\unsure{splitt|splitte} $\zeta$} annule $\zeta$.

% archive sheet number: 175
